\documentclass[11pt,a4paper]{article}

% ------------------------------------------------------------
% Packages
% ------------------------------------------------------------
\usepackage[margin=1in]{geometry}
\usepackage{amsmath,amssymb,amsthm,mathtools}
\usepackage{enumitem}
\usepackage{booktabs}
\usepackage{longtable}
\usepackage{array}
\usepackage{hyperref}
\usepackage{xcolor}
\usepackage{titlesec}
\usepackage{fancyhdr}

% ------------------------------------------------------------
% Hyperlinks
% ------------------------------------------------------------
\hypersetup{
    colorlinks=true,
    linkcolor=blue,
    urlcolor=blue,
    citecolor=blue
}

% ------------------------------------------------------------
% Page style
% ------------------------------------------------------------
\pagestyle{fancy}
\fancyhf{}
\lhead{Number Theory I}
\rhead{Course Structure}
\cfoot{\thepage}

% ------------------------------------------------------------
% Theorem environments
% ------------------------------------------------------------
\newtheorem{theorem}{Theorem}[section]
\newtheorem{definition}[theorem]{Definition}
\newtheorem{proposition}[theorem]{Proposition}
\newtheorem{lemma}[theorem]{Lemma}
\newtheorem{corollary}[theorem]{Corollary}

% ------------------------------------------------------------
% Custom commands
% ------------------------------------------------------------
\newcommand{\Z}{\mathbb Z}
\newcommand{\Q}{\mathbb Q}
\newcommand{\R}{\mathbb R}
\newcommand{\C}{\mathbb C}
\newcommand{\Fp}{\mathbb F_p}
\newcommand{\Fq}{\mathbb F_q}
\newcommand{\OK}{\mathcal O_K}
\newcommand{\OF}{\mathcal O_F}
\newcommand{\Gal}{\operatorname{Gal}}
\newcommand{\Tr}{\operatorname{Tr}}
\newcommand{\Norm}{\operatorname{N}}
\newcommand{\Spec}{\operatorname{Spec}}
\newcommand{\disc}{\operatorname{disc}}
\newcommand{\Cl}{\operatorname{Cl}}
\newcommand{\Hom}{\operatorname{Hom}}
\newcommand{\Aut}{\operatorname{Aut}}
\newcommand{\ord}{\operatorname{ord}}
\newcommand{\legendre}[2]{\left(\frac{#1}{#2}\right)}

% ------------------------------------------------------------
% Document
% ------------------------------------------------------------
\begin{document}

\begin{titlepage}
    \centering
    \vspace*{1.5cm}

    {\Huge\bfseries Number Theory I\par}
    \vspace{0.5cm}

    {\Large Full Course Structure, Modules, Concepts,\\
    Learning Objectives, Prerequisites, and Difficult Topics\par}

    \vspace{1.5cm}

    \rule{\textwidth}{0.5pt}

    \vspace{0.5cm}

    {\large
    A synthesized course structure based on the supplied
    tables of contents
    }

    \vspace{0.5cm}
    \rule{\textwidth}{0.5pt}

    \vfill

    \begin{minipage}{0.9\textwidth}
    \small
    \textbf{Scope.}
    The course is organized around the major themes appearing across
    the supplied sources: algebraic integers and Dedekind domains,
    completions and valuations, local fields, different and
    discriminant, cyclotomic fields, Minkowski theory, ideals,
    adeles and ideles, zeta functions and \(L\)-series, and
    local/global class field theory.  Advanced analytic and
    class-field-theoretic topics are included as the upper level
    of the course.
    \end{minipage}

    \vfill

    {\large \today\par}
\end{titlepage}

\tableofcontents
\newpage

% ============================================================
\section{Course Overview}
% ============================================================

\subsection{Course Description}

Number Theory I is a graduate-level course in modern algebraic and
analytic number theory.  The course develops the arithmetic of
integers and number fields and then extends the theory to local
fields, global fields, zeta functions, \(L\)-series, and class field
theory.

The conceptual progression is

\[
\boxed{
\text{Integers}
\longrightarrow
\text{Algebraic Integers}
\longrightarrow
\text{Ideals}
\longrightarrow
\text{Valuations}
\longrightarrow
\text{Local Fields}
\longrightarrow
\text{Global Fields}
\longrightarrow
\text{Zeta/\(L\)-Functions}
\longrightarrow
\text{Class Field Theory}
}
\]

The course should emphasize the interaction among three viewpoints:

\begin{enumerate}[label=\arabic*.]
    \item \textbf{Algebraic:} rings, ideals, field extensions,
    Galois groups, modules, and class groups;
    \item \textbf{Local:} valuations, completions, \(p\)-adic fields,
    ramification, local norms, and local reciprocity;
    \item \textbf{Analytic:} Dirichlet series, Euler products,
    zeta functions, \(L\)-series, functional equations, and
    density theorems.
\end{enumerate}

% ============================================================
\section{Source-Derived Course Architecture}
% ============================================================

The supplied tables of contents fall naturally into the following
large-scale blocks.

\begin{longtable}{p{0.18\textwidth}p{0.27\textwidth}p{0.45\textwidth}}
\toprule
\textbf{Block} & \textbf{Central Theme} & \textbf{Major Topics} \\
\midrule
I & Algebraic foundations &
Algebraic integers, localization, integral closure, prime ideals,
Dedekind domains, ideals, lattices, class numbers, units \\

II & Local arithmetic &
\(p\)-adic numbers, absolute values, valuations, completions,
local fields, Henselian fields, ramification \\

III & Algebraic extensions &
Different, discriminant, traces, norms, splitting of primes,
cyclotomic fields, ramification theory \\

IV & Geometry of numbers &
Product formula, lattices, parallelotopes, Minkowski theory,
ideal counting, class groups \\

V & Global arithmetic &
Adeles, ideles, ideal classes, global fields, function fields \\

VI & Analytic number theory &
Zeta functions, Dirichlet series, \(L\)-series, theta series,
Euler products, density of primes \\

VII & Class field theory &
Norm indices, Artin symbols, reciprocity laws, local and global
class fields \\

VIII & Advanced analytic theory &
Poisson summation, functional equations, Tauberian theorems,
Brauer--Siegel theorem, explicit formulas
\\
\bottomrule
\end{longtable}

% ============================================================
\section{Prerequisite Knowledge}
% ============================================================

\subsection{Essential Prerequisites}

Before beginning the course, students should be comfortable with:

\begin{enumerate}
    \item \textbf{Abstract algebra}
    \begin{itemize}
        \item groups, subgroups, quotient groups;
        \item rings, ideals, quotient rings;
        \item integral domains and fields;
        \item polynomial rings;
        \item field extensions;
        \item finite extensions;
        \item minimal and characteristic polynomials;
        \item basic Galois theory.
    \end{itemize}

    \item \textbf{Linear algebra}
    \begin{itemize}
        \item vector spaces and bases;
        \item linear transformations;
        \item determinants;
        \item matrices;
        \item dual spaces;
        \item finite-dimensional vector spaces.
    \end{itemize}

    \item \textbf{Elementary number theory}
    \begin{itemize}
        \item divisibility;
        \item Euclidean algorithm;
        \item greatest common divisors;
        \item congruences;
        \item prime factorization;
        \item arithmetic functions;
        \item Chinese remainder theorem;
        \item quadratic residues;
        \item basic finite fields.
    \end{itemize}

    \item \textbf{Basic analysis}
    \begin{itemize}
        \item sequences and series;
        \item convergence;
        \item continuity;
        \item basic integration;
        \item infinite products;
        \item complex numbers;
        \item elementary complex analysis.
    \end{itemize}
\end{enumerate}

\subsection{Recommended Prerequisites}

For the advanced modules, students should also know:

\begin{itemize}
    \item tensor products;
    \item modules over principal ideal domains;
    \item exact sequences;
    \item topological groups;
    \item Haar measure;
    \item Fourier analysis;
    \item basic complex analysis;
    \item infinite Galois theory.
\end{itemize}

\subsection{Prerequisite Dependency Structure}

\[
\begin{array}{c}
\text{Abstract Algebra}
\\
\downarrow
\\
\text{Algebraic Integers}
\rightarrow
\text{Ideals}
\rightarrow
\text{Dedekind Domains}
\\
\downarrow
\\
\text{Valuations}
\rightarrow
\text{Completions}
\rightarrow
\text{Local Fields}
\\
\downarrow
\\
\text{Different and Discriminant}
\rightarrow
\text{Ramification}
\\
\downarrow
\\
\text{Cyclotomic and Number Fields}
\\
\downarrow
\\
\text{Minkowski Theory}
\rightarrow
\text{Class Groups}
\\
\downarrow
\\
\text{Adeles and Ideles}
\\
\downarrow
\\
\text{Zeta Functions and \(L\)-Series}
\\
\downarrow
\\
\text{Class Field Theory}
\end{array}
\]

% ============================================================
\section{Major Course Learning Outcomes}
% ============================================================

By the end of the course, a successful student should be able to:

\begin{enumerate}
    \item Explain the arithmetic structure of algebraic number fields.
    \item Determine whether an element is integral over a base ring.
    \item Work with rings of algebraic integers.
    \item Factor ideals in Dedekind domains.
    \item Use localization to study arithmetic locally.
    \item Construct and work with \(p\)-adic numbers.
    \item Define and manipulate absolute values and valuations.
    \item Construct completions of valued fields.
    \item Explain the structure of local fields.
    \item Analyze unramified, tamely ramified, and more general
    field extensions.
    \item Compute traces, norms, differents, and discriminants.
    \item Study splitting and ramification of primes.
    \item Apply Minkowski's theorem to algebraic number fields.
    \item Understand the structure and significance of class groups
    and class numbers.
    \item Analyze units using Dirichlet's unit theorem.
    \item Work with cyclotomic fields and Gauss sums.
    \item Construct and use adeles and ideles.
    \item Understand the Dedekind zeta function and Dirichlet
    \(L\)-series.
    \item Explain Euler products and analytic continuation at an
    appropriate level.
    \item Understand the role of functional equations.
    \item Explain the Artin symbol and reciprocity maps.
    \item State and interpret local and global reciprocity laws.
    \item Explain the basic philosophy of class field theory.
    \item Connect algebraic, local, and analytic methods in number
    theory.
\end{enumerate}

% ============================================================
\section{Module I: Algebraic Integers}
% ============================================================

\subsection{Purpose}

This module establishes the algebraic foundation for arithmetic in
number fields.

\subsection{Core Concepts}

\begin{itemize}
    \item Gaussian integers \(\Z[i]\);
    \item algebraic integers;
    \item monic polynomials;
    \item integral dependence;
    \item integral closure;
    \item rings of integers \(\OK\);
    \item localization;
    \item prime and maximal ideals;
    \item quotient rings;
    \item Chinese remainder theorem;
    \item Galois extensions;
    \item Dedekind rings/domains;
    \item discrete valuation rings;
    \item factorization of prime ideals;
    \item projective modules;
    \item orders.
\end{itemize}

\subsection{Learning Objectives}

Students should be able to:

\begin{enumerate}
    \item Define an algebraic integer.
    \item Prove basic closure properties of algebraic integers.
    \item Determine \(\OK\) for basic number fields.
    \item Explain integral closure.
    \item Distinguish prime ideals from prime elements.
    \item Explain why unique factorization of elements may fail.
    \item Replace element factorization by ideal factorization.
    \item Describe Dedekind domains and their arithmetic significance.
    \item Work with localization at prime ideals.
\end{enumerate}

\subsection{Important Results}

\begin{itemize}
    \item Integral closure theorem;
    \item unique factorization of nonzero ideals in Dedekind domains;
    \item Chinese remainder theorem;
    \item basic structure theorem for finitely generated modules
    over suitable rings.
\end{itemize}

\subsection{Difficult Topics}

\begin{itemize}
    \item Integral closure;
    \item distinction between element and ideal factorization;
    \item projective modules over Dedekind domains;
    \item localization and its arithmetic interpretation.
\end{itemize}

% ============================================================
\section{Module II: Ideals, Lattices, Class Groups, and Units}
% ============================================================

\subsection{Core Concepts}

\begin{itemize}
    \item fractional ideals;
    \item ideal classes;
    \item class groups;
    \item class number;
    \item lattices;
    \item embeddings of number fields;
    \item Minkowski space;
    \item fundamental domains;
    \item covolume;
    \item Minkowski's theorem;
    \item ideal norms;
    \item counting ideals;
    \item Dirichlet's unit theorem.
\end{itemize}

\subsection{Learning Objectives}

Students should be able to:

\begin{enumerate}
    \item Define fractional ideals and ideal classes.
    \item Construct the ideal class group \(\Cl(K)\).
    \item Explain why the class group measures failure of
    principal ideal factorization.
    \item Embed a number field into a Euclidean space.
    \item Construct lattices associated with algebraic number fields.
    \item Apply Minkowski's convex body theorem.
    \item Obtain bounds on class groups and class numbers.
    \item State and interpret Dirichlet's unit theorem.
\end{enumerate}

\subsection{Difficult Topics}

\begin{itemize}
    \item Minkowski embeddings;
    \item volume computations;
    \item Minkowski's constant;
    \item relation between geometry and ideal theory;
    \item Dirichlet's unit theorem.
\end{itemize}

% ============================================================
\section{Module III: Completions and \(p\)-adic Numbers}
% ============================================================

\subsection{Core Concepts}

\begin{itemize}
    \item \(p\)-adic integers \(\Z_p\);
    \item \(p\)-adic numbers \(\Q_p\);
    \item \(p\)-adic absolute value;
    \item ultrametric inequality;
    \item Cauchy sequences;
    \item completions;
    \item discrete valuations;
    \item valuation rings;
    \item residue fields;
    \item filtrations;
    \item complete valued fields.
\end{itemize}

\subsection{Learning Objectives}

Students should be able to:

\begin{enumerate}
    \item Define the \(p\)-adic absolute value.
    \item Prove the ultrametric inequality.
    \item Construct \(\Q_p\) as a completion of \(\Q\).
    \item Describe \(\Z_p\).
    \item Work with \(p\)-adic expansions.
    \item Define a discrete valuation.
    \item Relate absolute values, valuations, and valuation rings.
    \item Explain the arithmetic meaning of completion.
\end{enumerate}

\subsection{Difficult Topics}

\begin{itemize}
    \item Non-Archimedean topology;
    \item convergence in \(\Q_p\);
    \item completions;
    \item equivalence between valuation and absolute-value
    formulations.
\end{itemize}

% ============================================================
\section{Module IV: Local Fields and Extensions}
% ============================================================

\subsection{Core Concepts}

\begin{itemize}
    \item local fields;
    \item finite extensions of \(\Q_p\);
    \item residue field extensions;
    \item unramified extensions;
    \item tamely ramified extensions;
    \item Henselian fields;
    \item extensions of valuations;
    \item Galois theory of valued fields;
    \item higher ramification groups.
\end{itemize}

\subsection{Learning Objectives}

Students should be able to:

\begin{enumerate}
    \item Characterize the basic examples of local fields.
    \item Extend valuations to algebraic extensions.
    \item Explain ramification index and residue degree.
    \item Distinguish unramified and ramified extensions.
    \item Describe tame ramification.
    \item Explain the role of Hensel's lemma.
    \item Relate local Galois groups to ramification.
    \item Describe higher ramification groups at a conceptual level.
\end{enumerate}

\subsection{Difficult Topics}

\begin{itemize}
    \item Henselian fields;
    \item extension of valuations;
    \item ramification groups;
    \item higher ramification theory;
    \item local Galois theory.
\end{itemize}

% ============================================================
\section{Module V: Trace, Norm, Different, and Discriminant}
% ============================================================

\subsection{Core Concepts}

\begin{itemize}
    \item field trace;
    \item field norm;
    \item trace pairing;
    \item complementary modules;
    \item inverse different;
    \item different ideal;
    \item discriminant;
    \item ramification;
    \item splitting of primes;
    \item separable extensions;
    \item inseparable extensions.
\end{itemize}

\subsection{Learning Objectives}

Students should be able to:

\begin{enumerate}
    \item Compute traces and norms of field elements.
    \item Relate trace and norm to multiplication operators.
    \item Define the different ideal.
    \item Compute discriminants in examples.
    \item Explain how the discriminant detects ramification.
    \item Relate local and global discriminants.
    \item Analyze the splitting behavior of primes.
\end{enumerate}

\subsection{Difficult Topics}

\begin{itemize}
    \item Different ideal;
    \item discriminant calculations;
    \item complementary modules;
    \item connection between discriminant and ramification;
    \item local/global compatibility.
\end{itemize}

% ============================================================
\section{Module VI: Cyclotomic Fields}
% ============================================================

\subsection{Core Concepts}

\begin{itemize}
    \item roots of unity;
    \item primitive \(n\)-th roots of unity;
    \item cyclotomic polynomials;
    \item cyclotomic fields;
    \item quadratic subfields;
    \item Galois groups;
    \item Gauss sums;
    \item ideal classes in cyclotomic fields.
\end{itemize}

\subsection{Learning Objectives}

Students should be able to:

\begin{enumerate}
    \item Define cyclotomic fields.
    \item Prove basic properties of cyclotomic polynomials.
    \item Determine the Galois group of a cyclotomic extension.
    \item Explain the appearance of quadratic fields.
    \item Understand the role of Gauss sums.
    \item Relate cyclotomic arithmetic to reciprocity laws.
\end{enumerate}

\subsection{Difficult Topics}

\begin{itemize}
    \item Gauss sums;
    \item ideal class relations;
    \item ramification in cyclotomic extensions.
\end{itemize}

% ============================================================
\section{Module VII: Global Fields, Function Fields, and Schemes}
% ============================================================

\subsection{Core Concepts}

\begin{itemize}
    \item number fields;
    \item function fields;
    \item places;
    \item completions;
    \item one-dimensional schemes;
    \item Dedekind rings;
    \item local-global viewpoints.
\end{itemize}

\subsection{Learning Objectives}

Students should be able to:

\begin{enumerate}
    \item Explain the analogy between number fields and function fields.
    \item Interpret places as local arithmetic data.
    \item Relate global fields to their completions.
    \item Recognize the role of one-dimensional schemes in arithmetic.
\end{enumerate}

\subsection{Difficult Topics}

\begin{itemize}
    \item Places and completions;
    \item function-field analogy;
    \item scheme-theoretic language.
\end{itemize}

% ============================================================
\section{Module VIII: Adeles and Ideles}
% ============================================================

\subsection{Core Concepts}

\begin{itemize}
    \item restricted direct products;
    \item local fields at places;
    \item adele rings;
    \item ideles;
    \item idele class groups;
    \item generalized ideal class groups;
    \item Galois actions on ideles.
\end{itemize}

\subsection{Learning Objectives}

Students should be able to:

\begin{enumerate}
    \item Explain why all completions of a global field are considered
    simultaneously.
    \item Define the adele ring.
    \item Define the idele group.
    \item Explain the difference between adeles and ideles.
    \item Relate ideal classes to idele classes.
    \item Describe the global arithmetic information encoded by ideles.
\end{enumerate}

\subsection{Difficult Topics}

This is one of the major conceptual transitions of the course.

\begin{itemize}
    \item Restricted direct products;
    \item topology on the adele ring;
    \item relation between ideals and ideles;
    \item Galois action on ideles;
    \item local-to-global synthesis.
\end{itemize}

% ============================================================
\section{Module IX: Zeta Functions and Dirichlet \(L\)-Series}
% ============================================================

\subsection{Core Concepts}

\begin{itemize}
    \item Dirichlet series;
    \item Euler products;
    \item Riemann zeta function;
    \item Dedekind zeta function;
    \item Dirichlet \(L\)-series;
    \item Hecke characters;
    \item Hecke \(L\)-series;
    \item Artin \(L\)-series;
    \item theta series;
    \item density of primes.
\end{itemize}

\subsection{Learning Objectives}

Students should be able to:

\begin{enumerate}
    \item Define Dirichlet series and Euler products.
    \item Derive the Euler product for the Riemann zeta function.
    \item Define the Dedekind zeta function.
    \item Explain how prime ideals enter zeta functions.
    \item Define Dirichlet characters and \(L\)-series.
    \item Explain the relationship between characters and \(L\)-functions.
    \item Understand the role of theta series.
    \item Explain the transition from Dirichlet \(L\)-series to Hecke
    and Artin \(L\)-series.
\end{enumerate}

\subsection{Difficult Topics}

\begin{itemize}
    \item Euler products;
    \item analytic properties of Dedekind zeta functions;
    \item Hecke characters;
    \item Artin \(L\)-series;
    \item density theorems.
\end{itemize}

% ============================================================
\section{Module X: Functional Equations and Analytic Methods}
% ============================================================

\subsection{Core Concepts}

\begin{itemize}
    \item Poisson summation;
    \item Fourier transforms;
    \item theta functions;
    \item functional equations;
    \item local additive duality;
    \item local multiplicative theory;
    \item global additive duality;
    \item Riemann--Roch connections;
    \item Tauberian theorems;
    \item non-vanishing of \(L\)-series.
\end{itemize}

\subsection{Learning Objectives}

Students should be able to:

\begin{enumerate}
    \item State and apply Poisson summation.
    \item Explain how Fourier analysis produces functional equations.
    \item Understand local and global Fourier duality.
    \item Derive the conceptual structure of a zeta-function
    functional equation.
    \item Explain the purpose of Tauberian theorems.
    \item Relate analytic behavior of \(L\)-functions to arithmetic
    information.
\end{enumerate}

\subsection{Difficult Topics}

\begin{itemize}
    \item Poisson summation;
    \item Fourier analysis on local fields;
    \item global additive duality;
    \item functional equations;
    \item Tauberian arguments.
\end{itemize}

% ============================================================
\section{Module XI: Analytic Theory of Prime Distribution}
% ============================================================

\subsection{Core Concepts}

\begin{itemize}
    \item Dirichlet integral;
    \item Tauberian theorem;
    \item non-vanishing of \(L\)-series;
    \item density of primes;
    \item prime-counting functions;
    \item asymptotic estimates;
    \item analytic continuation.
\end{itemize}

\subsection{Learning Objectives}

Students should be able to:

\begin{enumerate}
    \item Explain how analytic properties of zeta functions control
    prime distributions.
    \item Understand the role of non-vanishing results.
    \item Apply the conceptual form of a Tauberian theorem.
    \item Explain density statements for primes.
\end{enumerate}

% ============================================================
\section{Module XII: Brauer--Siegel Theory}
% ============================================================

\subsection{Core Concepts}

\begin{itemize}
    \item residue of Dedekind zeta functions;
    \item class number;
    \item regulator;
    \item discriminant;
    \item upper and lower residue estimates;
    \item normal extensions;
    \item Brauer--Siegel theorem.
\end{itemize}

\subsection{Learning Objectives}

Students should be able to:

\begin{enumerate}
    \item Explain the arithmetic quantities appearing in the
    residue of a Dedekind zeta function.
    \item Understand how class numbers and regulators enter
    analytic formulas.
    \item Explain the asymptotic relationship expressed by the
    Brauer--Siegel theorem.
\end{enumerate}

\subsection{Difficult Topics}

\begin{itemize}
    \item Analytic estimates;
    \item residues of zeta functions;
    \item comparison across field extensions;
    \item interaction of algebraic and analytic invariants.
\end{itemize}

% ============================================================
\section{Module XIII: Norm Indices and Reciprocity}
% ============================================================

\subsection{Core Concepts}

\begin{itemize}
    \item local norm groups;
    \item global norm indices;
    \item unit groups;
    \item exponential and logarithm functions;
    \item norm residue symbols;
    \item cyclic extensions;
    \item Artin symbol;
    \item reciprocity laws.
\end{itemize}

\subsection{Learning Objectives}

Students should be able to:

\begin{enumerate}
    \item Define local and global norm groups.
    \item Explain why norm indices are central to class field theory.
    \item Use local exponential and logarithm maps where appropriate.
    \item Understand the norm residue symbol.
    \item Define the Artin symbol conceptually.
    \item Explain the connection between Artin symbols and
    reciprocity laws.
\end{enumerate}

\subsection{Difficult Topics}

\begin{itemize}
    \item Local norm index;
    \item global norm index;
    \item norm residue symbols;
    \item Artin symbols;
    \item reciprocity.
\end{itemize}

% ============================================================
\section{Module XIV: Local Class Field Theory}
% ============================================================

\subsection{Core Concepts}

\begin{itemize}
    \item local reciprocity map;
    \item local norm groups;
    \item local Artin map;
    \item Hilbert symbol;
    \item norm residue symbol over \(\Q_p\);
    \item formal groups;
    \item generalized cyclotomic theory;
    \item local ramification theory.
\end{itemize}

\subsection{Learning Objectives}

Students should be able to:

\begin{enumerate}
    \item State the local reciprocity law.
    \item Describe the relationship between abelian extensions
    and norm groups.
    \item Explain the local Artin map.
    \item Interpret the Hilbert symbol.
    \item Relate local reciprocity to ramification.
    \item Understand the role of formal groups in local arithmetic.
\end{enumerate}

\subsection{Difficult Topics}

\begin{itemize}
    \item Local reciprocity;
    \item Hilbert symbol;
    \item formal groups;
    \item higher ramification;
    \item generalized cyclotomic theory.
\end{itemize}

% ============================================================
\section{Module XV: Global Class Field Theory}
% ============================================================

\subsection{Core Concepts}

\begin{itemize}
    \item idele classes;
    \item canonical pairings;
    \item reciprocity maps;
    \item Hasse reciprocity;
    \item global norm groups;
    \item class fields;
    \item Hilbert class field;
    \item principal ideal theorem;
    \item Brauer group;
    \item Hilbert symbols;
    \item power residue symbols.
\end{itemize}

\subsection{Learning Objectives}

Students should be able to:

\begin{enumerate}
    \item Explain the global reciprocity map.
    \item State Hasse's reciprocity law.
    \item Describe the Hilbert class field.
    \item Explain the principal ideal theorem in the class-field
    theoretic setting.
    \item Relate local reciprocity to global reciprocity.
    \item State the central existence and reciprocity theorems
    of global class field theory.
    \item Explain the role of the Brauer group.
\end{enumerate}

\subsection{Difficult Topics}

This is the most abstract algebraic portion of the course.

\begin{itemize}
    \item Idele class groups;
    \item canonical pairings;
    \item global reciprocity;
    \item existence theorem;
    \item Hilbert class fields;
    \item Brauer groups;
    \item compatibility of local and global reciprocity.
\end{itemize}

% ============================================================
\section{Module XVI: Advanced \(L\)-Functions}
% ============================================================

\subsection{Core Concepts}

\begin{itemize}
    \item proper abelian \(L\)-series;
    \item Artin \(L\)-series;
    \item induced characters;
    \item Hecke characters;
    \item Artin conductors;
    \item functional equations of Artin \(L\)-series;
    \item density theorems.
\end{itemize}

\subsection{Learning Objectives}

Students should be able to:

\begin{enumerate}
    \item Explain why \(L\)-functions are attached to algebraic
    and Galois-theoretic data.
    \item Define the Artin \(L\)-function conceptually.
    \item Explain induced characters and their contribution to
    \(L\)-series.
    \item Understand the meaning of an Artin conductor.
    \item Explain the structure of functional equations for
    Artin \(L\)-series.
    \item Relate \(L\)-functions to prime-density results.
\end{enumerate}

% ============================================================
\section{Module XVII: Explicit Formulas}
% ============================================================

\subsection{Core Concepts}

\begin{itemize}
    \item Weierstrass factorization;
    \item logarithmic derivatives;
    \item explicit formulas;
    \item Weil formula;
    \item zeros and poles of \(L\)-functions;
    \item prime sums.
\end{itemize}

\subsection{Learning Objectives}

Students should be able to:

\begin{enumerate}
    \item Explain the role of zeros of \(L\)-functions.
    \item Understand logarithmic derivatives of zeta functions.
    \item Interpret explicit formulas connecting zeros with primes.
    \item Explain the conceptual significance of Weil-type formulas.
\end{enumerate}

\subsection{Difficult Topics}

\begin{itemize}
    \item Weierstrass products;
    \item logarithmic derivatives;
    \item explicit formulas;
    \item zero--prime duality.
\end{itemize}

% ============================================================
\section{Recommended Course Sequence}
% ============================================================

A practical semester structure is the following.

\begin{longtable}{p{0.10\textwidth}p{0.24\textwidth}p{0.54\textwidth}}
\toprule
\textbf{Weeks} & \textbf{Module} & \textbf{Main Content} \\
\midrule
1--3 &
I &
Algebraic integers, integral closure, localization,
prime ideals, Dedekind domains \\

4--6 &
II &
Ideals, lattices, Minkowski theory, class groups,
units \\

7--9 &
III &
\(p\)-adic numbers, absolute values, valuations,
completions \\

10--12 &
IV &
Local fields, Henselian fields, unramified and
tamely ramified extensions \\

13--15 &
V &
Trace, norm, different, discriminant, splitting
and ramification \\

16--17 &
VI &
Cyclotomic fields and Gauss sums \\

18--19 &
VII &
Places, global fields, function fields \\

20--21 &
VIII &
Adeles, ideles, idele class groups \\

22--24 &
IX &
Zeta functions, Euler products, Dirichlet and
Hecke \(L\)-series \\

25--26 &
X &
Poisson summation and functional equations \\

27 &
XI &
Density of primes and Tauberian methods \\

28 &
XII &
Brauer--Siegel theory \\

29 &
XIII &
Norm indices, Artin symbols, reciprocity \\

30 &
XIV &
Local class field theory \\

31 &
XV &
Global class field theory \\

32 &
XVI--XVII &
Artin \(L\)-series and explicit formulas
\\
\bottomrule
\end{longtable}

\noindent
If the course is intended as a conventional one-semester
``Number Theory I'', Modules XIV--XVII should normally be treated
as an introduction or omitted in favor of deeper work on Modules
I--XIII.

% ============================================================
\section{Major Theorems and Results to Master}
% ============================================================

The following theorem list gives the mathematical backbone of the
course.

\subsection{Algebraic Number Theory}

\begin{enumerate}
    \item Chinese Remainder Theorem.
    \item Integral Closure Theorem.
    \item Unique ideal factorization in Dedekind domains.
    \item Structure theory of ideals in Dedekind domains.
    \item Dirichlet's Unit Theorem.
    \item Minkowski's Convex Body Theorem.
    \item Finiteness of the ideal class group.
\end{enumerate}

\subsection{Local Theory}

\begin{enumerate}
    \item Completion theorem.
    \item Hensel's Lemma.
    \item Structure theory of local fields.
    \item Classification of unramified extensions.
    \item Basic theory of tame ramification.
    \item Ramification theory for Galois extensions.
\end{enumerate}

\subsection{Field Extensions}

\begin{enumerate}
    \item Trace and norm formulas.
    \item Different--discriminant relationship.
    \item Prime decomposition in extensions.
    \item Ramification criteria.
    \item Basic cyclotomic field results.
\end{enumerate}

\subsection{Analytic Number Theory}

\begin{enumerate}
    \item Euler product for \(\zeta(s)\).
    \item Analytic properties of zeta functions.
    \item Functional equation.
    \item Poisson summation formula.
    \item Dirichlet \(L\)-series theory.
    \item Density theorems.
    \item Tauberian theorem applications.
    \item Brauer--Siegel theorem.
\end{enumerate}

\subsection{Class Field Theory}

\begin{enumerate}
    \item Norm-index theorems.
    \item Artin reciprocity.
    \item Local reciprocity law.
    \item Hilbert symbol theory.
    \item Hasse reciprocity.
    \item Existence theorem.
    \item Principal ideal theorem.
    \item Global reciprocity law.
\end{enumerate}

% ============================================================
\section{Difficult Topics: Difficulty Map}
% ============================================================

\begin{longtable}{p{0.30\textwidth}p{0.15\textwidth}p{0.43\textwidth}}
\toprule
\textbf{Topic} & \textbf{Difficulty} & \textbf{Reason} \\
\midrule
Algebraic integers & Moderate &
Requires integration of ring and field-extension theory \\

Dedekind domains & High &
Requires ideals, localization, modules, and factorization \\

Minkowski theory & High &
Combines geometry, linear algebra, and algebraic number theory \\

Dirichlet unit theorem & High &
Requires lattice methods and multiplicative structure \\

\(p\)-adic numbers & Moderate &
Requires adaptation from Archimedean to ultrametric analysis \\

Valuations & High &
Requires algebraic and topological viewpoints simultaneously \\

Henselian fields & High &
Combines valuations, polynomial factorization, and completeness \\

Ramification theory & Very High &
Requires local fields, Galois theory, valuations, and discriminants \\

Different and discriminant & High &
Abstract definitions must be connected to explicit computations \\

Adeles & Very High &
Requires simultaneous understanding of all completions \\

Ideles & Very High &
Requires groups, topology, valuations, and global arithmetic \\

Dedekind zeta functions & High &
Combines ideal theory with complex analysis \\

Hecke \(L\)-series & Very High &
Requires characters, ideals, harmonic analysis, and complex analysis \\

Poisson summation & High &
Requires Fourier analysis and careful analytic estimates \\

Functional equations & Very High &
Combines local and global analytic methods \\

Tauberian theory & Very High &
Requires sophisticated analytic estimates \\

Artin symbols & Very High &
Requires Galois theory and local/global norm theory \\

Local class field theory & Extremely High &
Combines local fields, Galois groups, norms, and reciprocity \\

Global class field theory & Extremely High &
Combines virtually every preceding part of the course \\

Artin \(L\)-series & Extremely High &
Requires Galois representations, characters, and analytic theory \\
\bottomrule
\end{longtable}

% ============================================================
\section{Conceptual Connections}
% ============================================================

\subsection{Element Factorization and Ideal Factorization}

A central motivation of the course is the replacement

\[
\text{factorization of elements}
\quad\longrightarrow\quad
\text{factorization of ideals}.
\]

In a number field \(K\), unique factorization of elements may fail,
while nonzero ideals in \(\OK\) have unique factorization.

This leads to

\[
\text{ideals}
\rightarrow
\text{ideal classes}
\rightarrow
\text{class group}
\rightarrow
\text{class number}.
\]

\subsection{Global and Local Arithmetic}

For a global field \(K\), each place \(v\) gives a completion

\[
K_v.
\]

Thus global arithmetic can be studied through the collection

\[
\{K_v\}_v.
\]

The course therefore develops the principle

\[
\boxed{
\text{Global problem}
\longrightarrow
\text{Local problems}
\longrightarrow
\text{Local solutions}
\longrightarrow
\text{Global synthesis}.
}
\]

\subsection{Valuations and Prime Ideals}

Valuations provide a local measure of divisibility:

\[
v_{\mathfrak p}(x).
\]

The corresponding valuation ring is

\[
\mathcal O_v
=
\{x : v(x)\geq 0\}.
\]

This creates the bridge

\[
\text{prime ideals}
\leftrightarrow
\text{valuations}
\leftrightarrow
\text{local rings}
\leftrightarrow
\text{completions}.
\]

\subsection{Discriminants and Ramification}

The discriminant provides a global measure of arithmetic
complexity, while ramification describes how primes behave locally.

Schematically,

\[
\disc(K)
\longleftrightarrow
\text{different}
\longleftrightarrow
\text{ramification}.
\]

\subsection{Zeta Functions and Arithmetic}

The Dedekind zeta function is

\[
\zeta_K(s)
=
\sum_{\mathfrak a\neq 0}
\frac{1}{N(\mathfrak a)^s}
=
\prod_{\mathfrak p}
\left(1-N(\mathfrak p)^{-s}\right)^{-1}.
\]

Thus ideals and prime ideals become analytic objects.

The progression is

\[
\text{ideals}
\rightarrow
\text{Euler product}
\rightarrow
\zeta_K(s)
\rightarrow
\text{analytic properties}
\rightarrow
\text{arithmetic information}.
\]

\subsection{Reciprocity}

The culmination of the course is the reciprocity philosophy:

\[
\boxed{
\text{Galois extensions}
\longleftrightarrow
\text{norm groups}
\longleftrightarrow
\text{idele classes}
\longleftrightarrow
\text{reciprocity maps}.
}
\]

% ============================================================
\section{Suggested Assessment Structure}
% ============================================================

\begin{longtable}{p{0.30\textwidth}p{0.20\textwidth}p{0.40\textwidth}}
\toprule
\textbf{Assessment} & \textbf{Suggested Weight} &
\textbf{Purpose} \\
\midrule
Weekly problem sets & 25\% &
Proofs, computations, and conceptual exercises \\

Midterm I & 15\% &
Algebraic number theory and ideal theory \\

Midterm II & 15\% &
Valuations, local fields, and ramification \\

Project / seminar & 10\% &
Independent study of a selected topic \\

Final examination & 25\% &
Comprehensive theoretical and computational mastery \\

Participation / presentations & 10\% &
Mathematical communication and discussion \\
\bottomrule
\end{longtable}

% ============================================================
\section{Problem-Solving Skills}
% ============================================================

Students should progressively acquire the following skills:

\begin{enumerate}
    \item \textbf{Computational skill:}
    calculate norms, traces, discriminants, ideal factorizations,
    valuations, and class groups in examples.

    \item \textbf{Proof skill:}
    prove structural statements concerning rings, fields, ideals,
    valuations, and Galois groups.

    \item \textbf{Local reasoning:}
    reduce a global arithmetic question to questions at individual
    primes or places.

    \item \textbf{Global reasoning:}
    reconstruct global information from local data.

    \item \textbf{Analytic reasoning:}
    translate arithmetic information into Dirichlet series,
    Euler products, and \(L\)-functions.

    \item \textbf{Structural reasoning:}
    recognize when a problem is naturally expressed using ideals,
    valuations, adeles, ideles, or Galois groups.

    \item \textbf{Research-level abstraction:}
    understand the conceptual role of reciprocity and class
    field theory.
\end{enumerate}

% ============================================================
\section{Core Topics Versus Advanced Topics}
% ============================================================

\subsection{Essential Core for Number Theory I}

The following topics should constitute the non-negotiable core:

\begin{enumerate}
    \item Algebraic integers.
    \item Integral closure.
    \item Prime ideals.
    \item Dedekind domains.
    \item Fractional ideals.
    \item Class groups and class numbers.
    \item Lattices and Minkowski theory.
    \item Dirichlet's unit theorem.
    \item \(p\)-adic numbers.
    \item Valuations.
    \item Completions.
    \item Local fields.
    \item Trace and norm.
    \item Different and discriminant.
    \item Prime decomposition and ramification.
    \item Cyclotomic fields.
    \item Basic zeta functions and \(L\)-series.
\end{enumerate}

\subsection{Advanced / Enrichment Topics}

These topics can be included toward the end of Number Theory I
or reserved for Number Theory II:

\begin{enumerate}
    \item Adeles and ideles.
    \item Abstract class field theory.
    \item Local class field theory.
    \item Global class field theory.
    \item Brauer groups.
    \item Formal groups.
    \item Artin \(L\)-series.
    \item Artin conductors.
    \item Tauberian theory.
    \item Brauer--Siegel theorem.
    \item Explicit formulas.
\end{enumerate}

% ============================================================
\section{Suggested Number Theory I / Number Theory II Boundary}
% ============================================================

\subsection{Number Theory I}

A coherent first course can end approximately at

\[
\boxed{
\text{Algebraic Number Theory}
+
\text{Local Theory}
+
\text{Basic Analytic Number Theory}.
}
\]

This includes

\begin{itemize}
    \item algebraic integers;
    \item ideals and Dedekind domains;
    \item class groups;
    \item Minkowski theory;
    \item units;
    \item valuations;
    \item \(p\)-adic numbers;
    \item local fields;
    \item discriminants;
    \item ramification;
    \item cyclotomic fields;
    \item basic zeta and \(L\)-functions.
\end{itemize}

\subsection{Number Theory II}

A natural continuation is

\[
\boxed{
\text{Adeles/Ideles}
+
\text{Advanced Analytic Theory}
+
\text{Class Field Theory}.
}
\]

This includes

\begin{itemize}
    \item adeles;
    \item ideles;
    \item Poisson summation;
    \item Tate-style local/global methods;
    \item Artin \(L\)-series;
    \item reciprocity laws;
    \item local class field theory;
    \item global class field theory;
    \item Brauer groups;
    \item explicit formulas.
\end{itemize}

% ============================================================
\section{Final Course Map}
% ============================================================

The entire curriculum can be summarized as follows:

\[
\begin{array}{cccccc}
\text{Algebra} &
\longrightarrow &
\text{Algebraic Integers} &
\longrightarrow &
\text{Ideals} &
\longrightarrow
\\[2mm]
&&
\downarrow &&
\downarrow
\\[2mm]
&&
\text{Dedekind Domains} &
\longrightarrow &
\text{Class Groups}
\\[2mm]
&&
\downarrow &&
\downarrow
\\[2mm]
\text{Valuations} &
\longrightarrow &
\text{Completions} &
\longrightarrow &
\text{Local Fields}
\\[2mm]
&&
\downarrow &&
\downarrow
\\[2mm]
&&
\text{Different/Discriminant} &
\longrightarrow &
\text{Ramification}
\\[2mm]
&&
\multicolumn{3}{c}{
\downarrow
}
\\[2mm]
&&
\text{Cyclotomic and Number Fields}
\\[2mm]
&&
\downarrow
\\[2mm]
&&
\text{Minkowski Theory}
\\[2mm]
&&
\downarrow
\\[2mm]
&&
\text{Adeles and Ideles}
\\[2mm]
&&
\downarrow
\\[2mm]
&&
\text{Zeta Functions and \(L\)-Series}
\\[2mm]
&&
\downarrow
\\[2mm]
&&
\text{Functional Equations}
\\[2mm]
&&
\downarrow
\\[2mm]
&&
\text{Reciprocity}
\\[2mm]
&&
\downarrow
\\[2mm]
&&
\text{Local Class Field Theory}
\\[2mm]
&&
\downarrow
\\[2mm]
&&
\boxed{\text{Global Class Field Theory}}
\end{array}
\]

% ============================================================
\section{Concluding Learning Goal}
% ============================================================

The ultimate goal of Number Theory I is not merely to accumulate
individual theorems.  Students should learn to recognize the
structural unity of arithmetic.

A typical arithmetic problem may be viewed through several
equivalent lenses:

\[
\boxed{
\begin{array}{c}
\text{integer / algebraic integer}
\\
\Updownarrow
\\
\text{ideal}
\\
\Updownarrow
\\
\text{valuation}
\\
\Updownarrow
\\
\text{local field}
\\
\Updownarrow
\\
\text{Galois group}
\\
\Updownarrow
\\
\text{\(L\)-function}
\\
\Updownarrow
\\
\text{reciprocity law}.
\end{array}
}
\]

The student completing the course should therefore be capable of
moving between algebraic, local, global, and analytic descriptions
of the same arithmetic phenomenon.

\end{document}